\documentclass[11pt,a4paper]{article}
\usepackage{fontspec,amsmath,amssymb,amsthm,microtype}
\setmainfont{lmroman10-regular.otf}[BoldFont=lmroman10-bold.otf,ItalicFont=lmroman10-italic.otf,BoldItalicFont=lmroman10-bolditalic.otf]
\usepackage{xcolor}
\usepackage[most]{tcolorbox}
\usepackage{mathrsfs}
\providecommand{\lt}{<}
\providecommand{\gt}{>}
\setmonofont{DejaVuSansMono.ttf}[BoldFont=DejaVuSansMono-Bold.ttf,ItalicFont=DejaVuSansMono-Oblique.ttf,BoldItalicFont=DejaVuSansMono-BoldOblique.ttf,Scale=0.86]
\newfontfamily\leanfallback{STIXTwoMath-Regular.otf}[Scale=0.86]
\definecolor{leanrule}{gray}{0.62}
\definecolor{leanback}{gray}{0.962}
\definecolor{leantitle}{gray}{0.22}
\definecolor{leanmuted}{gray}{0.38}
\newtcolorbox{leanbox}[1]{enhanced,breakable,colback=leanback,colframe=leanrule,
  boxrule=0.4pt,arc=0.6mm,left=2.4mm,right=2.4mm,top=1.8mm,bottom=1.8mm,
  before skip=0.8em,after skip=1em,fonttitle=\small\bfseries,coltitle=leantitle,
  colbacktitle=leanback,titlerule=0.3pt,title={#1}}
\newlength{\leancharwidth}
\newenvironment{leancode}{\par\smallskip\begingroup\ttfamily\footnotesize\raggedright
  \settowidth{\leancharwidth}{M}\setlength{\parskip}{0pt}\setlength{\parindent}{0pt}
  \frenchspacing}{\par\endgroup\smallskip}
\newcommand{\leanline}[2]{\hangindent=\dimexpr#1\leancharwidth+4\leancharwidth\relax
  \hangafter=1\noindent\hspace*{#1\leancharwidth}#2\par}
\newcommand{\leansource}[1]{\par{\footnotesize\color{leanmuted}#1}\par}
\newcommand{\leanlabel}[1]{\par\smallskip\noindent{\small\bfseries #1}\par}
\usepackage[margin=25mm]{geometry}
\usepackage[unicode,hidelinks]{hyperref}
\hypersetup{pdfauthor={LangevinLean Project}}
\setlength{\emergencystretch}{3em}
\setlength{\parskip}{0.4em}
\setlength{\parindent}{0pt}
\allowdisplaybreaks[2]
\raggedbottom
\title{\texorpdfstring{The random averaged Hessian on square-integrable vectors}{The random averaged Hessian on square-integrable vectors}}
\author{LangevinLean Project}
\date{}
\begin{document}\maketitle
Draft mathematical exposition; these notes may contain mistakes.


For actual random states \(X,Y\in L^2(Q;H)\), the field \(\omega\mapsto\overline A(X(\omega),Y(\omega))\) is measurable and bounded by \(\beta\). Multiplication by this field defines a bounded symmetric operator on \(L^2(Q;H)\), with the original curvature constants.



Earlier results used here: \cite[Section~3.13]{LangevinLean2026ULASC}.

\section{\texorpdfstring{Constructing multiplication by the actual field}{Constructing multiplication by the actual field}}


\textbf{Construction.} Set \((\mathcal A U)(\omega)=\overline A(X(\omega),Y(\omega))U(\omega)\). Joint continuity of the averaged Hessian gives measurable coefficients; its uniform operator bound gives an essentially bounded field. The estimate \(\|\mathcal A U(\omega)\|\le\beta\|U(\omega)\|\) proves square integrability and a bounded linear action on equivalence classes.

\begin{leanbox}{Lean: \texttt{averagedHessianField}, \texttt{averagedHessianField\_\allowbreak{}memLp}, \texttt{randomCurvature} and 1 more}
\leansource{\texttt{averagedHessianField} \quad WassersteinMixing/\allowbreak{}CurvatureLp.lean, lines 26–29}
\begin{leancode}
\leanline{0}{def\ averagedHessianField\ (X\ Y\ :\ Lp\ H\ 2\ Q)\ (ω\ :\ Ω)\ :\ H\ →L[ℝ]\ H\ :=}
\leanline{2}{averagedHessian\ V\ (X\ ω)\ (Y\ ω)}
\leanline{0}{\strut}
\leanline{0}{include\ P\ hV\ in}
\end{leancode}
\leansource{\texttt{averagedHessianField\_\allowbreak{}memLp} \quad WassersteinMixing/\allowbreak{}CurvatureLp.lean, lines 31–36}
\begin{leancode}
\leanline{0}{theorem\ averagedHessianField\_memLp\ (X\ Y\ :\ Lp\ H\ 2\ Q)\ :}
\leanline{4}{MemLp\ (averagedHessianField\ Q\ V\ X\ Y)\ ⊤\ Q}
\end{leancode}
\leansource{\texttt{randomCurvature} \quad WassersteinMixing/\allowbreak{}CurvatureLp.lean, lines 39–42}
\begin{leancode}
\leanline{0}{def\ randomCurvature\ (X\ Y\ :\ Lp\ H\ 2\ Q)\ :\ Lp\ H\ 2\ Q\ →L[ℝ]\ Lp\ H\ 2\ Q\ :=}
\leanline{2}{(ContinuousLinearMap.id\ ℝ\ (H\ →L[ℝ]\ H)).holderL\ Q\ ⊤\ 2\ 2}
\leanline{4}{((averagedHessianField\_memLp\ Q\ P\ V\ hV\ X\ Y).toLp}
\leanline{6}{(averagedHessianField\ Q\ V\ X\ Y))}
\end{leancode}
\leansource{\texttt{randomCurvature\_\allowbreak{}coe} \quad WassersteinMixing/\allowbreak{}CurvatureLp.lean, lines 45–53}
\begin{leancode}
\leanline{0}{theorem\ randomCurvature\_coe\ (X\ Y\ U\ :\ Lp\ H\ 2\ Q)\ :}
\leanline{4}{randomCurvature\ Q\ P\ V\ hV\ X\ Y\ U\ =ᵐ[Q]\ fun\ ω\ ↦}
\leanline{6}{averagedHessian\ V\ (X\ ω)\ (Y\ ω)\ (U\ ω)}
\end{leancode}
\leanlabel{How the Lean proof follows the argument}
\texttt{averaged\allowbreak{}Hessian\allowbreak{}Field\_\allowbreak{}mem\allowbreak{}Lp} establishes membership at exponent infinity. \texttt{random\allowbreak{}Curvature} uses bounded bilinear evaluation and the Hölder multiplication construction. \texttt{random\allowbreak{}Curvature\_\allowbreak{}coe} proves the exact almost-everywhere pointwise representative.
\end{leanbox}

\section{\texorpdfstring{Integrated forms, exact secants and adaptation}{Integrated forms, exact secants and adaptation}}


\textbf{Claim.} The constructed operator satisfies \(\|\mathcal A\|\le\beta\), is symmetric, and obeys

\[\alpha\|U\|_2^2\le\langle U,\mathcal A U\rangle_{L^2}\le\beta\|U\|_2^2,\qquad \mathcal A(X-Y)=\nabla V(X)-\nabla V(Y).\]

\textbf{Proof.} Integrate the pointwise norm and quadratic-form estimates using the exact representative. Pointwise Hessian symmetry gives symmetry of the integral pairing; the secant identity holds almost everywhere. If \(X,Y,U\) are measurable with respect to a sub-sigma-algebra, continuity of the field and evaluation makes their product measurable there too.

\begin{leanbox}{Lean: \texttt{randomCurvature\_\allowbreak{}norm\_\allowbreak{}le}, \texttt{randomCurvature\_\allowbreak{}symmetric}, \texttt{randomCurvature\_\allowbreak{}lower} and 3 more}
\leansource{\texttt{randomCurvature\_\allowbreak{}norm\_\allowbreak{}le} \quad WassersteinMixing/\allowbreak{}CurvatureLp.lean, lines 56–62}
\begin{leancode}
\leanline{0}{theorem\ randomCurvature\_norm\_le\ (X\ Y\ U\ :\ Lp\ H\ 2\ Q)\ :}
\leanline{4}{‖randomCurvature\ Q\ P\ V\ hV\ X\ Y\ U‖\ ≤\ P.beta\ *\ ‖U‖}
\end{leancode}
\leansource{\texttt{randomCurvature\_\allowbreak{}symmetric} \quad WassersteinMixing/\allowbreak{}CurvatureLp.lean, lines 65–73}
\begin{leancode}
\leanline{0}{theorem\ randomCurvature\_symmetric\ (X\ Y\ U\ W\ :\ Lp\ H\ 2\ Q)\ :}
\leanline{4}{inner\ ℝ\ (randomCurvature\ Q\ P\ V\ hV\ X\ Y\ U)\ W\ =}
\leanline{6}{inner\ ℝ\ U\ (randomCurvature\ Q\ P\ V\ hV\ X\ Y\ W)}
\end{leancode}
\leansource{\texttt{randomCurvature\_\allowbreak{}lower} \quad WassersteinMixing/\allowbreak{}CurvatureLp.lean, lines 76–86}
\begin{leancode}
\leanline{0}{theorem\ randomCurvature\_lower\ (X\ Y\ U\ :\ Lp\ H\ 2\ Q)\ :}
\leanline{4}{P.alpha\ *\ ‖U‖\ \textasciicircum{}\ 2\ ≤\ inner\ ℝ\ U\ (randomCurvature\ Q\ P\ V\ hV\ X\ Y\ U)}
\end{leancode}
\leansource{\texttt{randomCurvature\_\allowbreak{}upper} \quad WassersteinMixing/\allowbreak{}CurvatureLp.lean, lines 89–93}
\begin{leancode}
\leanline{0}{theorem\ randomCurvature\_upper\ (X\ Y\ U\ :\ Lp\ H\ 2\ Q)\ :}
\leanline{4}{inner\ ℝ\ U\ (randomCurvature\ Q\ P\ V\ hV\ X\ Y\ U)\ ≤\ P.beta\ *\ ‖U‖\ \textasciicircum{}\ 2}
\end{leancode}
\leansource{\texttt{randomCurvature\_\allowbreak{}secant} \quad WassersteinMixing/\allowbreak{}CurvatureLp.lean, lines 96–109}
\begin{leancode}
\leanline{0}{theorem\ randomCurvature\_secant\ [IsFiniteMeasure\ Q]\ (X\ Y\ :\ Lp\ H\ 2\ Q)\ :}
\leanline{4}{randomCurvature\ Q\ P\ V\ hV\ X\ Y\ (X\ -\ Y)\ =}
\leanline{6}{ULA.lpDrift\ Q\ (gradient\ V)\ hV.gradientLipschitz\ X\ -}
\leanline{8}{ULA.lpDrift\ Q\ (gradient\ V)\ hV.gradientLipschitz\ Y}
\end{leancode}
\leansource{\texttt{randomCurvature\_\allowbreak{}adapted} \quad WassersteinMixing/\allowbreak{}CurvatureLp.lean, lines 113–120}
\begin{leancode}
\leanline{0}{theorem\ randomCurvature\_adapted\ (m\ :\ MeasurableSpace\ Ω)\ (X\ Y\ U\ :\ Lp\ H\ 2\ Q)}
\leanline{4}{(hX\ :\ AEStronglyMeasurable[m]\ X\ Q)\ (hY\ :\ AEStronglyMeasurable[m]\ Y\ Q)}
\leanline{4}{(hU\ :\ AEStronglyMeasurable[m]\ U\ Q)\ :}
\leanline{4}{AEStronglyMeasurable[m]\ (randomCurvature\ (m0\ :=\ m0)\ Q\ P\ V\ hV\ X\ Y\ U)\ Q}
\end{leancode}
\leanlabel{How the Lean proof follows the argument}
The norm and pairing lemmas use the pointwise representative and the integral expression for the real square norm. \texttt{random\allowbreak{}Curvature\_\allowbreak{}secant} invokes the actual segment secant identity. \texttt{random\allowbreak{}Curvature\_\allowbreak{}adapted} proves measurability under the supplied sub-sigma-algebra, without an adaptation field on the operator.
\end{leanbox}


Earlier results used here: \cite[Section~3.13.2]{LangevinLean2026ULASC}; \cite[Section~3.13.3]{LangevinLean2026ULASC}.

\begin{thebibliography}{ULA-LSI}
\bibitem[ULA-SC]{LangevinLean2026ULASC}
LangevinLean Project. \emph{Unadjusted Langevin under Strong Convexity}. LL-01, 2026. Draft report, content revision d34ead06987c. \url{https://odakyildiz.github.io/LangevinLean/docs/pdf/UnadjustedLangevin/FullProof.pdf}.
\end{thebibliography}

\end{document}
